\documentclass[]{ConspectBook}
\let\enquote\emph

%  1) Charges
%  2) ChargeSystem
%  3) EBMotion
%  4) ElEnergy
%  5) ElMagInduction
%  6) EMOscillations
%  7) EMWaves
%  8) FileldInDielectrics
%  9) FileldInMetall
% 10) MagEnergy
% 11) MagFieldMedia
% 12) MagFieldVacuum
% 13) MaxwellEqns
% 14) Radiation
% 15) Relativity
% 16) SteadyCurrent
% 17) SteadyEfield
% 18) Telegraph
% 19) Waveguides

\newcommand{\includechapter}[1]{\input{#1/#1}}
\newcommand{\includetikz}[2]{\input{#1/tikz/#2.tikz}}


\begin{document}

%\includechapter{Tensions}
\section{Трифазний струм}

Трифазну систему змінного струму із шістьма проводами винайшов у 1888 році Тесла. У 1889 році російський винахідник М. Доливо-Добровольський запропонував трифазну систему з трьома проводами. На електротехнічній виставці у Франкфурті-на-Майні в 1891 році він продемонстрував дослідну високовольтну електропередачу трифазного струму протяжністю 175 км. Для цього він використав трифазний генератор, який мав потужність 230 кВт при напрузі 95 В.

Будову трифазного синхронного генератора схематично показано на рис.~\ref{pic:ThreePhaseGenerator}. У пазах осердя статора розташовані три однакові обмотки I, II і III, зсунуті одна відносно одної на \(120^\circ\). При рівномірній частоті обертання ротора в обмотках статора індукуються гармонічні ЕРС однакової частоти, зсунуті за фазою відносно одна одної на кут \(2\pi/3\). Кожна з ЕРС знаходиться у своїй фазі періодичного процесу, тому частоту називають просто «фазою». Також «фазами» називають обмотки генератора і приймача. Кінці обмоток (фаз) підключені к трифазній електричній мережі — сукупності трьох однофазних електричних кіл. Фази з'єднують зіркою або трикутником.

\begin{figure}[h!]\centering
\begin{tikzpicture}[>=Latex, thick]

  % ---------- Параметри геометрії (міняючи їх, перемальовується вся картинка) ----------
  \def\Rlam{3.0}   % зовнішній радіус осердя статора (тонке коло)
  \def\Rout{2.85}  % зовнішній радіус обмоткового шару
  \def\Rin{2.15}   % внутрішній радіус обмоткового шару (межа повітряного зазору)
  \def\gapI{90}    % кут "розриву" I  (фаза 1)
  \def\gapII{330}  % кут "розриву" II (фаза 2)
  \def\gapIII{210} % кут "розриву" III (фаза 3)
  \def\halfgap{6}  % половина кутової ширини розриву, град.

  % ---------- Осердя статора ----------
  \draw (0,0) circle (\Rlam);

  % ---------- Кільце обмотки (заливка через even-odd) ----------
  \draw[even odd rule, fill=cyan!55!blue!25] (0,0) circle (\Rout) (0,0) circle (\Rin);

  % ---------- Петлі обмотки: три групи по 120°, кожна намальована циклом ----------
  \foreach \a in {96,102,...,204}
    {\draw[cyan!40!black] ({(\Rout+0.10)*cos(\a)},{(\Rout+0.10)*sin(\a)}) circle (0.11);}
  \foreach \a in {216,222,...,324}
    {\draw[cyan!40!black] ({(\Rout+0.10)*cos(\a)},{(\Rout+0.10)*sin(\a)}) circle (0.11);}
  \foreach \a in {336,342,...,444}
    {\draw[cyan!40!black] ({(\Rout+0.10)*cos(\a)},{(\Rout+0.10)*sin(\a)}) circle (0.11);}

  % ---------- Ротор: двополюсна "арахісоподібна" форма r(t)=a+b*cos(2(t-90)) ----------
  \draw[fill=magenta!45]
    plot[domain=0:360,samples=180,smooth,variable=\t]
    ({(1.15+0.85*cos(2*(\t-90)))*cos(\t)},{(1.15+0.85*cos(2*(\t-90)))*sin(\t)}) -- cycle;

  % отвір під вал
  \draw[fill=white] (0,0) circle (0.18);

  % три штрихові осі, співвісні з розривами обмотки
  \foreach \a in {\gapI,\gapII,\gapIII}
    {\draw[dashed] ({0.18*cos(\a)},{0.18*sin(\a)}) -- ({2.05*cos(\a)},{2.05*sin(\a)});}

  \node at (0,0.55) {\Large С};
  \node at (0,-0.55) {\Large Ю};

  % ---------- Позначки 120° (дуга + повернутий підпис) ----------
  \draw[<->] ({2.5*cos(96+10)},{2.5*sin(96+10)}) arc (96+10:204-10:2.5);
  \node[rotate=60]  at ({3.35*cos(150)},{3.35*sin(150)}) {$120^\circ$};

  \draw[<->] ({2.5*cos(216+10)},{2.5*sin(216+10)}) arc (216+10:324-10:2.5);
  \node[rotate=0]   at ({3.35*cos(270)},{3.35*sin(270)}) {$120^\circ$};

  \draw[<->] ({2.5*cos(336+10)},{2.5*sin(336+10)}) arc (336+10:444-10:2.5);
  \node[rotate=-60] at ({3.35*cos(30)},{3.35*sin(30)}) {$120^\circ$};

  % ---------- Стрілки I, II, III біля розривів обмотки ----------
  \draw[->] ({3.6*cos(\gapI)},{3.6*sin(\gapI)})   -- ({(\Rlam+0.05)*cos(\gapI)},{(\Rlam+0.05)*sin(\gapI)});
  \node at ({3.85*cos(\gapI)},{3.85*sin(\gapI)}) {I};

  \draw[->] ({3.6*cos(\gapII)},{3.6*sin(\gapII)})  -- ({(\Rlam+0.05)*cos(\gapII)},{(\Rlam+0.05)*sin(\gapII)});
  \node at ({3.85*cos(\gapII)},{3.85*sin(\gapII)}) {II};

  \draw[->] ({3.6*cos(\gapIII)},{3.6*sin(\gapIII)}) -- ({(\Rlam+0.05)*cos(\gapIII)},{(\Rlam+0.05)*sin(\gapIII)});
  \node at ({3.85*cos(\gapIII)},{3.85*sin(\gapIII)}) {III};

  % ================= Виводи до навантажень =================

  % ---- Фаза 1 -> Z1 (найкоротший, "внутрішній" контур) ----
  \draw ({\Rlam*cos(96)},{\Rlam*sin(96)}) -- ({\Rlam*cos(96)},2.75) -- (6.2,2.75);
  \node at (1.0,2.95) {$I_1$};
  \draw ({\Rlam*cos(84)},{\Rlam*sin(84)}) -- ({\Rlam*cos(84)},2.15) -- (6.2,2.15);
  \node at (1.0,2.35) {$I_1$};
  \draw (6.2,1.85) rectangle (6.8,3.05);
  \node at (6.5,2.45) {$Z_1$};

  % ---- Фаза 2 -> Z2 (середній контур) ----
  \draw ({\Rlam*cos(336)},{\Rlam*sin(336)}) -- ({\Rlam*cos(336)},0.05) -- (5.0,0.05);
  \node at (2.3,0.25) {$I_2$};
  \draw ({\Rlam*cos(324)},{\Rlam*sin(324)}) -- ({\Rlam*cos(324)},-0.75) -- (5.0,-0.75);
  \node at (2.3,-0.95) {$I_2$};
  \draw (5.0,-1.05) rectangle (5.6,0.35);
  \node at (5.3,-0.35) {$Z_2$};

  % ---- Фаза 3 -> Z3 (найдовший контур, огинає всю картинку зверху і знизу) ----
  \draw ({\Rlam*cos(204)},{\Rlam*sin(204)}) -- ({\Rlam*cos(204)},3.9) -- (7.4,3.9) -- (7.4,-1.0);
  \node at (-2.1,4.1) {$I_3$};
  \draw ({\Rlam*cos(216)},{\Rlam*sin(216)}) -- ({\Rlam*cos(216)},-3.9) -- (7.4,-3.9) -- (7.4,-2.6);
  \node at (-2.1,-4.1) {$I_3$};
  \draw (7.1,-2.6) rectangle (7.7,-1.0);
  \node at (7.4,-1.8) {$Z_3$};

\end{tikzpicture}
	\caption{Схематичне зображення трифазного синхронного генератора\label{pic:ThreePhaseGenerator}}
\end{figure}

Зіркою називається таке з'єднання, при якому кінці обмоток генератора і приймача з'єднують в одну спільну точку, яку називають нейтральною точкою або \emph{нейтраллю}. Кінці обмоток трифазного приймача (електродвигуна, трансформатора и т.д.) з'єднують також в одну точку — нейтраль (рис.~\ref{pic:ThreePhaseStar}). Провід, що з'єднує два нейтралі, називається \emph{нейтральним}. Три проводи, що з'єднують обмотки генератора і приймача, називаються \emph{лінійними проводами}.

\begin{figure}[h!]\centering
	\begin{circuitikz}[>=latex, scale=1.2]
		% Генератор (зірка)
		\coordinate (N1) at (0,0);
		\draw (N1) node[above] {\(O\)} -- (120:1.5) node[above right] {\(I\)};
		\draw (N1) -- (240:1.5) node[below right] {\(II\)};
		\draw (N1) -- (0:1.5) node[right] {\(III\)};
		% Лінійні проводи
		\draw (120:1.5) -- ++(3,0) coordinate (L1);
		\draw (240:1.5) -- ++(3,0) coordinate (L2);
		\draw (0:1.5) -- ++(3,0) coordinate (L3);
		% Нейтральний провід
		\draw (N1) -- ++(0,-2.5) coordinate (N2);
		% Приймач (зірка)
		\coordinate (N3) at (6,-2.5);
		\draw (N2) -- (N3) node[right] {\(O'\)};
		\draw (L1) -- ++(0,-1) to[R=$Z_1$] ++(0,-1.5) -- (N3);
		\draw (L2) -- ++(0,-1) to[R=$Z_2$] ++(0,-1.5) -- (N3);
		\draw (L3) -- ++(0,-1) to[R=$Z_3$] ++(0,-1.5) -- (N3);
		% Підписи
		\node at (3,1.5) {\(I_{\text{л}1}\)};
		\node at (3,0.5) {\(I_{\text{л}2}\)};
		\node at (3,-0.5) {\(I_{\text{л}3}\)};
		\node at (1.5,-1.5) {\(I_\text{H}\)};
		% Напруги
		\draw[<->] (L1) ++(1,0) -- ++(0,-1.5) node[midway, right] {\(U_\text{ф}\)};
		\draw[<->] (L1) ++(1.5,0.2) -- ++(0,-2) node[midway, right] {\(U_\text{л}\)};
	\end{circuitikz}
	\caption{З'єднання зіркою\label{pic:ThreePhaseStar}}
\end{figure}

Напруга \(U_\text{ф}\) між лінійним проводом і нейтраллю називається \emph{фазною}. Напруга \(U_\text{л}\) між двома лінійними проводами називається \emph{лінійною}. На векторній діаграмі, зображеній на рис.~\ref{pic:ThreePhaseVectorDiagram}, показані три вектори \(\mathbf{U}_\text{ф}\), повернуті один відносно одного на кут \(120^\circ\), і вектор \(\mathbf{U}_\text{л}\). З діаграми видно, що \(U_\text{л} = \sqrt{3} \cdot U_\text{ф}\). Зокрема, якщо ефективна напруга \(U_\text{ф} = 220\) В, то \(U_\text{л} = 380\) В.

\begin{figure}[h!]\centering
	\begin{tikzpicture}[>=latex, scale=1.5]
		% Центр
		\coordinate (O) at (0,0);
		% Вектори фазних напруг
		\draw[->, thick, blue] (O) -- (90:1.5) node[above] {\(\mathbf{U}_{\text{ф}1}\)};
		\draw[->, thick, blue] (O) -- (210:1.5) node[below left] {\(\mathbf{U}_{\text{ф}2}\)};
		\draw[->, thick, blue] (O) -- (330:1.5) node[below right] {\(\mathbf{U}_{\text{ф}3}\)};
		% Вектор лінійної напруги (між 1 та 2)
		\draw[->, thick, red] (90:1.5) -- (210:1.5) node[midway, left] {\(\mathbf{U}_\text{л}\)};
		% Кути
		\draw (0.4,0) arc (0:120:0.4);
		\node at (60:0.6) {\small \(120^\circ\)};
		\draw (0.4,0) arc (0:-120:0.4);
		\node at (-60:0.6) {\small \(120^\circ\)};
		% Позначення точок
		\node at (90:1.7) {\(1\)};
		\node at (210:1.7) {\(2\)};
		\node at (330:1.7) {\(3\)};
	\end{tikzpicture}
	\caption{Векторна діаграма напруг\label{pic:ThreePhaseVectorDiagram}}
\end{figure}

Аналогічно струм \(I_\text{ф}\), що тече в обмотці генератора або приймача, називається фазним струмом, а струм \(I_\text{л}\) в лінійному проводі — лінійним струмом. При з'єднанні зіркою \(I_\text{л} = I_\text{ф}\), а вектор струму в нейтралі \(\mathbf{I}_\text{H} = \mathbf{I}_{\text{л}1} + \mathbf{I}_{\text{л}2} + \mathbf{I}_{\text{л}3}\).

Якщо опори \(Z_1\), \(Z_2\) і \(Z_3\) приймача рівні між собою, то таке навантаження називають \emph{симетричним}. В цьому випадку струм \(\mathbf{I}_\text{H}\) в нейтралі відсутній, і нейтральний провід можна видалити.

При з'єднанні генератора і приймача зіркою (рис.~\ref{pic:ThreePhaseStar}) використовується три проводи. Фазне і лінійне напруження збігаються (\(U_\text{ф} = U_\text{л}\)), а лінійний струм більший за фазний: \(I_\text{л} = \sqrt{3} \cdot I_\text{ф}\).

\begin{figure}[h!]\centering
	\begin{circuitikz}[>=latex, scale=1.2]
		% Генератор (трикутник)
		\coordinate (A) at (0,0);
		\coordinate (B) at (2,0);
		\coordinate (C) at (1,1.732);
		\draw (A) -- (B) -- (C) -- cycle;
		\node at (A) [below left] {\(I\)};
		\node at (B) [below right] {\(II\)};
		\node at (C) [above] {\(III\)};
		% Виводи
		\draw (A) -- ++(-1.5,0) coordinate (L1);
		\draw (B) -- ++(1.5,0) coordinate (L2);
		\draw (C) -- ++(0,1) coordinate (L3);
		% Приймач (трикутник)
		\coordinate (A') at (6,0);
		\coordinate (B') at (8,0);
		\coordinate (C') at (7,1.732);
		\draw (A') to[R=$Z_1$] (B');
		\draw (B') to[R=$Z_2$] (C');
		\draw (C') to[R=$Z_3$] (A');
		% З'єднання
		\draw (L1) -- (A');
		\draw (L2) -- (B');
		\draw (L3) -- (C');
		% Підписи
		\node at (3,0.5) {\(I_{\text{л}1}\)};
		\node at (3,-0.5) {\(I_{\text{л}2}\)};
		\node at (4,1.5) {\(I_{\text{л}3}\)};
		\draw[<->] (A) ++(0.5,0) -- ++(0,-1) node[midway, right] {\(U_\text{ф}\)};
		\draw[<->] (L1) ++(0,0.5) -- ++(1.5,0) node[midway, above] {\(U_\text{л}\)};
	\end{circuitikz}
	\caption{З'єднання трикутником\label{pic:ThreePhaseDelta}}
\end{figure}

Передача електричної енергії на дальні відстані по трифазних колах є вигіднішою, ніж передача енергії по однофазних колах. Трифазні синхронні генератори, трифазні синхронні та асинхронні двигуни і трансформатори (див. далі) простіші у виробництві, економічні та надійні в експлуатації. У трифазних системах досить просто отримати обертове магнітне поле, дія якого на провідники зі струмом покладена в основу принципу роботи асинхронних і синхронних електродвигунів.



\end{document}

